%% conference_review_en.tex — English conference version of the review
%% "Sensitive Indicators of Soil Health".
%% Source: paper/conference_review.tex (shortened Russian conference version).
%% Build: pdflatex conference_review_en.tex (twice).

\documentclass[english]{ai4science_ysda}
\usepackage{float}

\InformationInEnglish
\hyphenation{physicochemical}

\TITLE{Sensitive Indicators of Soil Health: Responses to Management and
Limitations}
\ShortTITLE{Sensitive indicators of soil health}
\AUTHOR{}

\Abstract{Soil health is described through physical, chemical and biological
properties, and the choice of indicators for monitoring remains debated. This
review compares three groups of indicators: physicochemical (soil organic
carbon and its normalised forms), labile carbon fractions and microbiological
indicators. For each group, the review lists the management
practices that produce a measurable response and the conditions that limit
interpretation. In the reviewed studies, labile fractions responded to reduced
tillage and organic matter input. Microbiological indicators responded to a
wider set of practices. Physicochemical indicators provided the basis for
assessment after normalisation to the pedoclimatic context. A comparison table
lists the practices, the direction of the response and the limitations.}

\KEYWORDS{soil health, indicators, soil organic carbon, labile carbon
fractions, microbiological indicators}

\begin{document}

\section*{Introduction and scope}

\porozhdeno{} Soil health is defined as the capacity of soil to function as a living system
and to deliver ecosystem services, so monitoring requires indicators that
respond to management and remain linked to long-term changes in soil condition.

Soil organic carbon (SOC) is among the most widely used indicators, together
with soil acidity (pH) and available phosphorus and potassium \cite{bunemann2018}.
The absolute SOC content depends on climate, soil type and land use, and this
dependence complicates the definition of target contents. This review compares
indicators from three groups and, for each group, lists the management practices
that produce a response and the conditions under which the interpretation fails.

\section{Preparation of the review}

\porozhdeno{} The review was prepared during the course ``AI assistants for
researchers''. The project developed an agent for literature work and for the
preparation of a review on soil health indicators. The full review is based on a
corpus of 47 scientific articles. The conference text is a shortened version of
the full review and keeps the indicator groups, their comparison, the conditions
of application and the limitations.

%% conference_abbreviations_en.tex — abbreviations used in the English
%% conference version of the review. Included with \input from
%% conference_review_en.tex.

\begin{table}[H]
  \centering
  \caption{Abbreviations used in the review.}
  \label{tab:abbreviations}
  \small
  \setlength{\tabcolsep}{2pt}
  \renewcommand{\arraystretch}{0.85}
  \begin{tabular}{@{}>{\raggedright\arraybackslash}p{0.15\textwidth}
      >{\raggedright\arraybackslash}p{0.28\textwidth}
      >{\raggedright\arraybackslash}p{0.15\textwidth}
      >{\raggedright\arraybackslash}p{0.28\textwidth}@{}}
    \hline
    \textbf{Abbreviation} & \textbf{Definition} & \textbf{Abbreviation} &
    \textbf{Definition} \\
    \hline
    SOC & soil organic carbon & STN & soil total nitrogen \\
    SOC/clay & soil organic carbon-to-clay ratio & O/T SOC &
    observed-to-typical soil organic carbon ratio \\
    DOC & dissolved organic carbon & Hy-DOC & hydrophilic dissolved organic
    carbon \\
    POXC & permanganate-oxidisable carbon & HWEC & hot-water-extractable
    carbon \\
    POMC & particulate organic matter carbon & MBC & microbial biomass carbon \\
    PLFA & phospholipid fatty acids & B-GLU & $\beta$-glucosidase \\
    FDA & fluorescein diacetate hydrolysis & URE & urease \\
    EOM & exogenous organic matter & RR & crop residue retention \\
    CC & cover crops & CR & crop rotation diversification \\
    TR & tillage reduction & IOS & innovative or organic systems \\
    LUCAS & Land Use/Cover Area Frame Survey & & \\
    \hline
  \end{tabular}
\end{table}


\section{Physicochemical indicators and SOC normalisation}

\porozhdeno{} Normalising soil organic carbon to clay content accounts for its
association with fine mineral particles, including clay and silt
\cite{feeney2024}. The soil organic
carbon-to-clay ratio (SOC/clay) serves as a threshold of carbon saturation.
Normalisation to clay alone is biased by texture and climate and can produce
misleading estimates of soil structure and carbon storage potential
\cite{feeney2024}.

The observed-to-typical soil organic carbon ratio (O/T SOC) is an alternative:
the ratio of the observed carbon content to the content typical for a
pedoclimatic zone \cite{feeney2024}. Typical contents were derived from the Land
Use/Cover Area Frame Survey (LUCAS) programme (2009--2018, 19\,855 mineral
topsoil samples of 0--20~cm). Twelve pedoclimatic zones were identified, with
typical SOC contents of 5.99--39.65~g~kg$^{-1}$ \cite{feeney2024}. The indicator
is less sensitive to clay content, land use and climate, and it resolves
differences in porosity and carbon stocks better \cite{feeney2024}.
\section{Labile carbon fractions}

\porozhdeno{} Management-induced changes are often difficult to detect in total
SOC content, and rapidly cycling pools are therefore measured
\cite{bongiorno2019}. Ten European long-term field experiments compared five
fractions: dissolved organic carbon (DOC), hydrophilic dissolved organic carbon
(Hy-DOC), permanganate-oxidisable carbon (POXC), hot-water-extractable carbon
(HWEC) and particulate organic matter carbon (POMC) \cite{bongiorno2019}.

Reduced tillage and high organic matter input increased the contents of labile
fractions relative to conventional tillage and low input. POXC and POMC were the
most responsive to both factors. POXC also correlated most closely with chemical,
physical and biological indicators \cite{bongiorno2019}. POXC measurement is
inexpensive and rapid.

The response depends on sampling conditions, and some studies detected no effect
of tillage and fertilisation on labile fractions. The chemical composition and
seasonal dynamics of POXC have received little study \cite{bongiorno2019}.

\section{Microbiological indicators}

\porozhdeno{} A meta-analysis of 166 articles assessed the response of active
organic matter pools and microbiological indicators to six practices: input of
exogenous organic matter (EOM), crop residue retention (RR), cover crops (CC),
crop rotation diversification (CR), tillage reduction (TR) and a transition to
innovative or organic systems (IOS) \cite{cusset2024}. Climatic factors showed
no significant effect on the response in this meta-analysis. Clay and organic
matter contents, pH, experiment duration and sampling depth modulated the
response \cite{cusset2024}.

Physicochemical indicators such as SOC and soil total nitrogen (STN), and
microbiological indicators, both responded to EOM, RR, CC and TR. Under CC, SOC
content increased by 7\,\%, STN by 8\,\%, microbial biomass carbon (MBC) by
14\,\%, total phospholipid fatty acids (PLFA) by 19\,\% and $\beta$-glucosidase
(B-GLU) activity by 11\,\% \cite{cusset2024}. CR and IOS did not affect SOC and
STN contents, but increased MBC by 7--34\,\%, PLFA by 8--24\,\% and B-GLU
activity by 58\,\% under the transition to IOS \cite{cusset2024}. These values
apply to the listed practices and to the conditions of the meta-analysis.

Dehydrogenase, phosphatase and $\beta$-glucosidase activities reflect the effects
of tillage, organic fertilisers and crop rotation. These assays are sensitive to sample preparation and storage
conditions \cite{bhupenchandra2026}. Functional genes of the carbon, nitrogen and
phosphorus cycles responded to fertiliser treatments in a ten-year field
experiment \cite{jia2025}, but their measurement requires molecular methods. The
composition of microbial communities responds to management, and the effect size
depends on the biome \cite{trivedi2016}. Priming effects complicate the
interpretation, because microbial biomass drives carbon and nitrogen turnover
\cite{kuzyakov2010}.

\section{Comparison and conditions of application}

\porozhdeno{} Physicochemical indicators provide the basis for assessing the
organic status of soil but respond slowly and require normalisation to the
pedoclimatic context \cite{feeney2024,cardoso2013}. Labile fractions respond to
reduced tillage and organic matter input \cite{bongiorno2019}. Microbiological
indicators respond to a wider set of practices, including those that leave SOC
and STN contents unchanged in short-term experiments \cite{cusset2024}. The
reviewed studies do not support a single ranking of indicators.

Agreement with the wider literature is partial. At one long-term site, reduced
tillage did not change organic matter content, cation exchange capacity or POXC
in a conventional system \cite{pearsons2023}. In a synthesis of long-term
experiments in Europe and China, organic matter input affected earthworm
abundance most strongly, followed by crop yield, organic matter content and
aggregate stability. The absence of tillage increased aggregate stability and
organic matter content in the topsoil, while crop yield decreased
\cite{bai2018}.

For large-scale monitoring, a minimum set of three indicators has been proposed:
SOC, aggregate stability and the 24-hour carbon mineralisation potential
\cite{bagnall2023}.

\begin{table}[!htbp]
  \centering
  \caption{Comparison of indicators: management practices, direction of the
  response and limitations. Practice codes are defined in Section~4; all practice
  and indicator codes are listed in Table~\ref{tab:abbreviations}.}
  \label{tab:indicators}
  \small
  \setlength{\tabcolsep}{2pt}
  \renewcommand{\arraystretch}{0.85}
  \begin{tabular}{@{}>{\raggedright\arraybackslash}p{0.16\textwidth}
      >{\raggedright\arraybackslash}p{0.21\textwidth}
      >{\raggedright\arraybackslash}p{0.19\textwidth}
      >{\raggedright\arraybackslash}p{0.20\textwidth}
      >{\raggedright\arraybackslash}p{0.12\textwidth}@{}}
    \hline
    \textbf{Indicator} & \textbf{Practice or condition} &
    \textbf{Direction of response} & \textbf{Limitation} &
    \textbf{Source} \\
    \hline
    SOC, STN & EOM, RR, CC, TR & contents increase; weak response to CR & slow
    response; climate and texture & \cite{cusset2024} \\
    SOC/clay & normalisation to clay & threshold of carbon saturation & bias
    with texture and climate & \cite{feeney2024} \\
    O/T SOC & pedoclimatic zones & separates porosity and carbon stocks & not
    for organic soils and subsoils & \cite{feeney2024} \\
    POXC & TR, EOM & increase under TR and high EOM & sampling conditions and
    timing & \cite{bongiorno2019} \\
    POMC & TR, EOM & among the most responsive fractions & heterogeneity of
    residues & \cite{bongiorno2019} \\
    DOC, HWEC & TR, EOM & weaker response than POXC and POMC & soil properties
    and season & \cite{bongiorno2019} \\
    MBC & EOM, RR, CC, CR, IOS & increase, also without SOC increase & cost and
    expertise & \cite{cusset2024} \\
    PLFA & EOM, RR, CC, CR, IOS & increase under CR and IOS & complex protocol &
    \cite{cusset2024} \\
    Enzymes (B-GLU, FDA, URE) & EOM, RR, CC, CR, IOS & increase; B-GLU highest
    under IOS & assay standardisation & \cite{cusset2024,bhupenchandra2026} \\
    Earthworms & EOM, no-tillage, grass & abundance increases & fewer data;
    pesticides & \cite{bai2018} \\
    Aggregate stability & no-tillage, organic farming & stability increases &
    depth and texture & \cite{bai2018} \\
    \hline
  \end{tabular}

  \vspace{2pt}
  {\footnotesize Practice and indicator codes are defined in
  Table~\ref{tab:abbreviations}.\par}
\end{table}

\section{Verification of statements}

\porozhdeno{} Scientific claims, numerical results, measurement units and
bibliographic details were checked manually against the original scientific
publications. The checks covered the reported management practices, study
conditions and limitations relevant to each cited result.

The shortened English version was compared with the full review to preserve the
meaning and scope of the source material. Automated checks verified the
structure of the full review and the format of its indicator table. The
conference document was compiled twice and checked for page count, unresolved
references and layout errors.

\section{Limitations}

\porozhdeno{} The completeness of the literature coverage was not assessed with a formal
systematic-search protocol, and comparability across studies is limited by
differences in soil and climate, experiment duration, management practices and
measurement methods. The O/T SOC indicator was developed for mineral topsoil.
Typical contents are not defined for organic soils and subsoil horizons, and the
classes depend on the accuracy of digital maps \cite{feeney2024}. The merged
LUCAS surveys exclude the analysis of change over time \cite{feeney2024}.

Labile fractions depend on sampling conditions and seasonal dynamics
\cite{bongiorno2019}, and microbiological indicators require more time, expertise
and cost \cite{cusset2024}. The meta-analysis used a fixed sampling depth
\cite{cusset2024}. The effects of microbial indicators depend on the biome
\cite{trivedi2016}. Functional genes require specialised infrastructure
\cite{jia2025}. The resolution of bioindicators limits their transfer between
conditions \cite{bhaduri2022}. Sensitivity estimates are qualitative, and no
single scale is introduced.

\section*{Conclusion}

\porozhdeno{} Labile carbon fractions and microbiological indicators provide
complementary information on soil responses to management. POXC and POMC
responded strongly to reduced tillage and organic matter inputs in the reviewed
experiments. Microbial biomass carbon, phospholipid fatty acids and enzyme
activities responded to a broader range of practices, including crop
diversification and transitions to innovative or organic systems, for which
responses in total soil organic carbon and nitrogen were weaker or absent.

Soil organic carbon remains a basic indicator of soil condition. Its
normalisation against pedoclimatic reference values provides context for
interpreting differences between sites. Combining these measurements with
labile carbon fractions and microbiological indicators links the assessment of
soil organic status with the detection of management responses. This combination
supports the selection of indicators suited to the practices and soil conditions
under assessment.

\begin{aiusage}[Tags: \porozhdeno{} --- fragment generated by a model,
\istochnik{} --- fragment taken from a source.]
  Abstract, Sections 1--7; translation into English and language editing &
    \porozhdeno{} & open-weight LLM via Yandex AI Studio; version not recorded &
    compared with the full review; claims, numbers and units checked manually
    against the original publications; no human authorship claimed \\
  Table~\ref{tab:abbreviations} and Table~\ref{tab:indicators} & \porozhdeno{} &
    same tool & from the table in Section~7 of the full review; abbreviations
    collected from the text \\
  Bibliography & \istochnik{} & \texttt{bibitem} entries transferred from the
    source \texttt{.tex} & bibliographic details checked manually against the
    original publications \\
  Compilation and checks & \porozhdeno{} & pdflatex; checks/check.py covers the
    full review and its indicator table & two compilations; page count and
    references in the report \\
\end{aiusage}

\begin{thebibliography}{13}
\bibitem{feeney2024}
Feeney, C.J., Bentley, L., De Rosa, D., Panagos, P., Emmett, B.A., Thomas, A.,
Robinson, D.A. (2024). Benchmarking soil organic carbon (SOC) concentration
provides more robust soil health assessment than the SOC/clay ratio at European
scale. \emph{Science of the Total Environment}, 951, 175642.

\bibitem{bai2018}
Bai, Z., Caspari, T., Ruiperez Gonzalez, M., Batjes, N.H., Mäder, P.,
Bünemann, E.K., et al. (2018). Effects of agricultural management practices on
soil quality: A review of long-term experiments for Europe and China.
\emph{Agriculture, Ecosystems and Environment}, 265, 1--7.

\bibitem{bongiorno2019}
Bongiorno, G., Bünemann, E.K., Oguejiofor, C.U., Meier, J., Gort, G., Comans, R.,
Mäder, P., Brussaard, L., de Goede, R. (2019). Sensitivity of labile carbon
fractions to tillage and organic matter management and their potential as
comprehensive soil quality indicators across pedoclimatic conditions in Europe.
\emph{Ecological Indicators}, 99, 38--50.

\bibitem{cusset2024}
Cusset, E., Bennegadi-Laurent, N., Recous, S., Bernard, P.-Y., Perrin, A.-S.,
Tscheiller, R., Trinsoutrot-Gattin, I., Riah-Anglet, W. (2024). Which soil
microbial indicators should be included in routine laboratory tests to support
the transition to sustainable management of arable farming systems?
A meta-analysis. \emph{Ecological Indicators}, 167, 112706.

\bibitem{bhaduri2022}
Bhaduri, D., Sihi, D., Bhowmik, A., Verma, B.C., Munda, S., Dari, B. (2022).
A review on effective soil health bio-indicators for ecosystem restoration and
sustainability. \emph{Frontiers in Microbiology}, 13, 938481.

\bibitem{cardoso2013}
Cardoso, E.J.B.N., Vasconcellos, R.L.F., Bini, D., Miyauchi, M.Y.H.,
Santos, C.A., Alves, P.R.L., de Paula, A.M., Nakatani, A.S., Pereira, J.M.,
Nogueira, M.A. (2013). Soil health: looking for suitable indicators. What
should be considered to assess the effects of use and management on soil
health? \emph{Scientia Agricola}, 70(4), 274--289.

\bibitem{jia2025}
Jia, J., de Goede, R., Li, Y., Zhang, J., Wang, G., Zhang, J., Creamer, R.
(2025). Unlocking soil health: Are microbial functional genes effective
indicators? \emph{Soil Biology and Biochemistry}, 204, 109768.

\bibitem{trivedi2016}
Trivedi, P., Delgado-Baquerizo, M., Anderson, I.C., Singh, B.K. (2016).
Response of soil properties and microbial communities to agriculture:
Implications for primary productivity and soil health indicators.
\emph{Frontiers in Plant Science}, 7, 990.

\bibitem{bunemann2018}
Bünemann, E.K., Bongiorno, G., Bai, Z., Creamer, R., De Deyn, G., de Goede, R.,
et al. (2018). Soil quality --- a critical review. \emph{Soil Biology and
Biochemistry}, 120, 105--125.

\bibitem{bagnall2023}
Bagnall, D.K., Rieke, E.L., Morgan, C.L.S., Liptzin, D.L., Cappellazzi, S.B.,
Honeycutt, C.W. (2023). A minimum suite of soil health indicators for North
American agriculture. \emph{Soil Security}, 10, 100084.

\bibitem{pearsons2023}
Pearsons, K.A., et al. (2023). A tale of two systems: Does reducing tillage
affect soil health differently depending on the system? \emph{Soil \& Tillage
Research}, 226, 105562.

\bibitem{kuzyakov2010}
Kuzyakov, Y. (2010). Priming effects: Interactions between living and dead
organic matter. \emph{Soil Biology and Biochemistry}, 42(9), 1363--1371.

\bibitem{bhupenchandra2026}
Bhupenchandra, I., Sharma, L.D., Salam, M.D., Kumar, S., Chongtham, S.K.,
Choudhary, A.K., Devi, H.L., Sinyorita, S., Dutta, P., Bouket, A.C.,
Minkina, T., Keswani, C. (2026). Unravelling the role of soil enzymatic
activity as a vital bioindicator of soil health and fertility in agricultural
systems. \emph{Applied Soil Ecology}, 219, 106816.
\end{thebibliography}

\end{document}
